\section{Simulation study}\label{sec:simulation}

The simulation study was designed to evaluate the predictive performance
of the proposed GG-KM model under different cure mechanisms, latent-count
distributions, and sample sizes. Specifically, we investigated three
questions. First, we examined whether a simpler predictor performs best
when its functional form is correctly specified. Second, we evaluated
whether GG-KM provides robust predictions when the true cure-probability
function is not adequately represented by the simpler models. Third, we
assessed the extent to which predictive performance depends on the
latent-count distribution and its degree of overdispersion. The
generalized gamma latency distribution was correctly specified and held
fixed across all scenarios. Therefore, differences in predictive
performance primarily reflect the ability of each predictor to recover
the covariate-dependent cure mechanism.

Each data-generating scenario was defined by one cure-probability surface,
one latent-count distribution, and one sample size. We considered three
cure-probability surfaces, linear, quadratic, and periodic. We also
considered three latent-count distributions, Poisson, negative binomial
with dispersion $\phi=0.25$, and negative binomial with dispersion
$\phi=1.00$. The two negative-binomial settings represent modest and
stronger overdispersion, respectively, whereas the Poisson distribution
represents the equidispersed limiting case. The sample sizes were
$n=300$ and $n=600$, resulting in $18$ data-generating scenarios.

In every scenario, two covariates were generated independently as
$x_1,x_2\sim\operatorname{Uniform}(-1,1)$. The bounded uniform domain
deliberately differs from the historical source because it gives all
three geometries a common domain for level-set and recovery comparisons.
The source used standard-normal covariates and a covariate-dependent
Weibull latency; retaining those choices would confound cure-predictor
shape with unmodelled latency heterogeneity.

The true cure probability was generated from one of three fixed surfaces.
Specifically, with $\operatorname{expit}(v)=\{1+\exp(-v)\}^{-1}$, the
designs were
\begin{equation}
	\begin{aligned}
		p_c^{\mathrm{lin}}(\bm{x})
		&=\operatorname{expit}(0.15-1.00x_1+0.75x_2),\\
		p_c^{\mathrm{quad}}(\bm{x})
		&=\operatorname{expit}(0.60-1.45x_1^2-1.05x_2^2+0.55x_1x_2),\\
		p_c^{\mathrm{per}}(\bm{x})
		&=\operatorname{expit}\{0.10+1.15\sin(\pi x_1)+0.80\cos(\pi x_2)\}.
	\end{aligned}
	\label{eq:simulation-cure-surfaces}
\end{equation}
For Poisson generation, the latent-count mean was
$\theta(\bm{x})=-\log p_c(\bm{x})$; for negative-binomial generation, it
was $\theta(\bm{x})=\{p_c(\bm{x})^{-\phi}-1\}/\phi$. This mapping
preserves the same cure-probability surface across count distributions,
so comparisons do not confound dispersion with a different cure target.

Figure~\ref{fig:cure-surfaces} shows the three surfaces on the common
covariate square. The linear surface increases along one tilted gradient
and has parallel straight level sets, so the log-linear predictor is
correctly specified for it. The quadratic surface has closed elliptical
level sets generated by the two squared terms and the interaction, and no
additive predictor reproduces that geometry. The periodic surface reverses
its gradient twice inside the domain and produces alternating regions of
high and low cure probability, a shape an additive spline can follow one
coordinate at a time but a degree-two map cannot. The three designs
therefore separate correct linear specification, smooth nonadditivity and
oscillation while holding the latency distribution fixed. They are not
matched on prevalence: the mean of $p_c(\bm{x})$ over the square is
0.533 for the linear surface, 0.448 for the quadratic surface and 0.520
for the periodic surface, with attainable ranges $[0.168,0.870]$,
$[0.079,0.646]$ and $[0.136,0.886]$. Scores are consequently comparable
between methods within a scenario and not between scenarios.

\GGKMWideFigure{cure-surfaces.pdf}{Frozen cure-probability surfaces used
in the simulation design. Panels show the linear, quadratic, and
periodic functions in Equation~\eqref{eq:simulation-cure-surfaces} over
$(x_1,x_2)\in[-1,1]^2$. All panels use the same $p_c(\bm{x})$ colour
scale; superimposed contour lines provide a colour-independent encoding
of the surface geometry. Contours are labelled at $p_c=0.25$, $0.50$ and
$0.60$, the same levels used in Figure~\ref{fig:cure-surface-recovery},
and each panel title reports the mean of $p_c(\bm{x})$ over the square.
}{fig:cure-surfaces}{Three contour panels show cure probability over two
covariates. The linear panel has diagonal straight contours, the quadratic
panel has curved contours with an interaction, and the periodic panel has
repeated peaks and troughs. Each panel title states the mean cure
probability for that surface.
}

All latent-cause activation times followed the generalized gamma
distribution with $(a,d,p)=(1.6,1.3,1.1)$. Independent censoring times
were exponential with mean 3.5, and observed time was the minimum of the
event and censoring times. The inverse population distribution was used
to generate finite event times directly for susceptible individuals, while
individuals in the limiting cure mass were assigned no event. Every
replicate retained the true individual cure probability and population
survival curve for evaluation.

Algorithm~\ref{alg:simulation-generation} states the implemented
inverse-transform generator. Sampling from the improper population
distribution avoids simulating a potentially large latent count while
preserving its cure mass and finite-event distribution exactly.
\begin{algorithm*}[tbp]
\caption{Right-censored simulation from a promotion-time population distribution}\label{alg:simulation-generation}
\begin{algorithmic}[1]
\Require Scenario, sample size $n$, generalized gamma parameters $(a,d,p)$, censoring scale, and deterministic replication seed
\For{$i=1,\ldots,n$}
  \State Draw $x_{i1},x_{i2}\stackrel{\mathrm{ind}}{\sim}\operatorname{Uniform}(-1,1)$ and evaluate the declared $p_c(\bm{x}_i)$
  \State Set $\theta_i=-\log p_c(\bm{x}_i)$ for Poisson counts, or $\theta_i=\{p_c(\bm{x}_i)^{-\phi}-1\}/\phi$ for negative-binomial counts
  \State Draw $U_i\sim\operatorname{Uniform}(0,1)$ from the population event-time distribution
  \If{$U_i\geq1-p_c(\bm{x}_i)$}
    \State Set $T_i=\infty$ (the draw falls in the cure mass)
  \Else
    \State Set $q_i=-\log(1-U_i)/\theta_i$ for Poisson counts
    \State Set $q_i=\{(1-U_i)^{-\phi}-1\}/(\phi\theta_i)$ for negative-binomial counts
    \State Set $T_i=F_{GG}^{-1}(q_i;a,d,p)$
  \EndIf
  \State Draw $C_i$ from the declared independent censoring distribution
  \State Record the observed time $t_i=\min(T_i,C_i)$, indicator $\delta_i=I(T_i\leq C_i)$, $p_c(\bm{x}_i)$, and the true population survival grid
\EndFor
\State Return the complete observed sample, generating truth, scenario identifier, and replication seed
\end{algorithmic}
\end{algorithm*}

The fitted comparison set shared the same generalized gamma latency and comprised
GG-KM predictors with linear, radial-basis, Laplacian, exponential, Cauchy,
polynomial, and wavelet kernels, each with Poisson and negative-binomial latent
counts, together with additive cubic-spline Poisson and negative-binomial models
and the classical mixture-cure and promotion-time comparators. Linear-kernel
candidates used $\lambda\in\{0.3,3\}$. Radial-basis, Laplacian, exponential, and
Cauchy candidates used $\gamma\in\{0.25,1\}$ with
$\lambda\in\{0.3,3\}$. Polynomial candidates additionally used
degree $q\in\{2,3\}$ with $c_0=1$, while wavelet candidates used
$A\in\{0.5,1,2\}$, in both cases with $\lambda\in\{0.3,3\}$.
Spline candidates used three or five training-fold quantile knots with
$\lambda\in\{0.3,3\}$. All tuned methods used the same two-fold inner tuning
criterion and evaluation-grid size. The promotion-time support-vector method of
Pal and Aselisewine was not included because the reproducibility audit found no
public executable implementation at the article or arXiv records, and a new
reimplementation would not provide a controlled software comparison.
\cite{pal2023}

Each generated sample was split once into a common event-stratified
$70\%$ training set and $30\%$ test set for all fitted methods. Tuning,
feature-map fitting, and preprocessing were confined to the training set,
and the IPCW integration horizon was set from that set at the smaller of
its 90th observed-time percentile and the limit at which estimated
censoring survival remained at least 0.05. Primary test-set estimands were
integrated Brier score, cure AUC, and cure-probability mean squared error.

A target-size diagnostic pilot used five planned replicates per pilot
scenario and was not pooled with the final experiment; at least four
replicates were successful in every pilot cell. The largest observed
between-replicate standard deviations were 0.0418 for integrated Brier
score and 0.0332 for cure-probability mean squared error. Requiring
projected Monte Carlo standard errors no larger than 0.003 and 0.005,
respectively, implied at least 195 replicates for integrated Brier score
and 45 for cure-probability mean squared error. The more restrictive
criterion therefore required 195 replicates, so the final count was fixed
at 250 replicates per scenario before final estimation. At
$R=250$, the corresponding projected Monte Carlo standard errors were
0.00265 for integrated Brier score and 0.00210 for cure-probability mean
squared error. For cure AUC, the largest pilot standard deviation was
0.1211, corresponding to a projected Monte Carlo standard error of
0.00766 at $R=250$. Final summaries report the realized Monte Carlo
standard error for every primary metric and paired differences computed
within each shared replication.

\subsection{Simulation results}\label{sec:simulation-results}

The primary final experiment contained 81,000 method rows from
250 replications in each of 18 scenarios. The largest realized
Monte Carlo standard errors were 0.00175 for the IPCW integrated
Brier score, 0.00434 for cure-probability RMSE, and
0.00694 for cure AUC. Of the 81,000 fits, 80,842 completed
successfully and 158 failed; among the successful fits, 80,832
satisfied the explicit convergence criterion and 10 did not.
Method-specific successful-fit totals ranged from 4,467 to 4,498
of 4,500. Mean per-fit runtime ranged from 1.13 to 40.32 seconds
across methods, with a maximum of 520.48 seconds.

Table~\ref{tab:simulation-best} reports the best observed mean IPCW
integrated Brier score, cure root mean squared error, and cure AUC in each
scenario. For the linear surface, a linear-kernel GG-KM attained the lowest
mean IPCW integrated Brier score and cure root mean squared error in all six
scenarios. For the quadratic surface, a polynomial-kernel GG-KM attained the
lowest cure root mean squared error in all six scenarios and the lowest
integrated Brier score in five of six scenarios, with an RBF-kernel GG-KM
attaining the lowest integrated Brier score in the remaining case. For the
periodic surface, an additive spline model attained the lowest integrated
Brier score and cure root mean squared error in all six scenarios. For cure
AUC, a polynomial-kernel GG-KM attained the highest mean in all six linear
scenarios and all six periodic scenarios. On the quadratic surface,
polynomial-kernel GG-KM models attained the highest mean cure AUC in four
of six scenarios, whereas wavelet-kernel GG-KM models attained the highest
mean in the remaining two.

\GGKMTable{sections/generated/simulation-best-table}

Figure~\ref{fig:simulation-paired-differences} shows paired differences
between the RBF GG-KM and the count-matched linear-kernel GG-KM,
additive-spline, and polynomial-kernel GG-KM comparators across the six
generating-law and sample-size cells. Each contrast is defined as the RBF
GG-KM value minus the comparator value; hence, negative differences favor
the RBF model for IPCW integrated Brier score and cure root mean squared
error, whereas positive differences favor it for cure AUC.

For the linear cure surface, the linear-kernel GG-KM consistently
outperformed the RBF model in integrated Brier score and cure root mean
squared error, while the spline and polynomial comparisons were generally
closer. Cure AUC also favored the linear and polynomial kernels, whereas
the RBF--spline differences remained comparatively small. Under the
quadratic surface, the RBF GG-KM substantially improved on the linear
kernel across all three metrics, but its differences from the spline and
polynomial models were much smaller; the polynomial kernel generally had
lower cure root mean squared error. For the periodic surface, the RBF
model improved markedly on the linear and polynomial kernels in integrated
Brier score and cure root mean squared error, particularly at $n=600$,
but remained inferior to the spline. Cure AUC showed a different pattern,
with the RBF model clearly outperforming the linear kernel and, at
$n=600$, the Poisson polynomial-kernel GG-KM, while its differences from
the spline remained comparatively small. Overall, the paired comparisons
show that relative performance depends strongly on both the geometry of
the cure surface and the evaluation metric.

\GGKMWideFigure
{simulation_rbf_paired_differences_with_auc.pdf}
{Count-matched paired predictive differences of the RBF GG-KM against the linear-kernel GG-KM, additive-spline, and polynomial-kernel GG-KM comparators. Columns correspond to the linear, quadratic, and periodic cure surfaces, while rows show IPCW integrated Brier score, cure root mean squared error, and cure AUC differences. Each contrast is defined as the RBF GG-KM value minus the corresponding comparator value, with Poisson and negative-binomial fits compared separately. Error bars span $\pm1.96$ paired Monte Carlo standard errors, and a horizontal line marks zero. Negative differences favor the RBF GG-KM for integrated Brier score and cure root mean squared error, whereas positive differences favor it for cure AUC.}
{fig:simulation-paired-differences}
{
A three-row by three-column grid displays count-matched paired
differences between the RBF GG-KM and linear-kernel GG-KM,
additive-spline, and polynomial-kernel GG-KM comparators. Columns
represent linear, quadratic, and periodic cure surfaces; rows represent
integrated Brier score, cure root mean squared error, and cure AUC.
Poisson and negative-binomial contrasts are distinguished by marker and
line type, with error bars spanning $\pm1.96$ paired Monte Carlo standard
errors around a horizontal zero line. Negative values favor the RBF model
for integrated Brier score and cure root mean squared error, whereas
positive values favor it for cure AUC.
}

Figure~\ref{fig:cure-surface-recovery} maps where in the
covariate square each predictor loses accuracy. Every panel
fills the $41\times41$ evaluation grid with the pointwise root
mean squared error of the predicted cure probability against
the generating surface, averaged over the valid Poisson
replications at $n=600$, on a scale shared by all twelve panels.
Solid lines repeat the true level sets at $p_c=0.25$, $0.50$
and $0.60$, and the dashed line is the mean estimated
$p_c=0.50$ level, which separates an error of magnitude from a
displacement of shape. Failed fits reduce the pointwise
$n_{\mathrm{valid}}$ rather than being removed from the total
number of replications. The panel means summarize the gridwise
recovery of the cure surface. On the linear surface the
linear-kernel GG-KM attains 0.0490, compared with 0.0671 for
the additive spline, 0.0694 for the polynomial-kernel GG-KM,
and 0.0730 for the RBF GG-KM. On the quadratic surface, the
RBF GG-KM attains the lowest mean error, 0.0802, followed
closely by the additive spline at 0.0816 and the
polynomial-kernel GG-KM at 0.0827, whereas the linear-kernel
GG-KM increases to 0.1282. On the periodic surface, the
additive spline performs best at 0.0824, followed by the RBF
GG-KM at 0.0959, while the linear- and polynomial-kernel GG-KM
predictors increase to 0.1504 and 0.1510, respectively. The
spatial pattern adds what these scalar summaries omit. A
misaligned predictor does not spread its error evenly: the
linear-kernel GG-KM shows concentrated error on the quadratic
and periodic surfaces, and its estimated $p_c=0.50$ level is
absent on the quadratic surface because the fitted cure surface
does not reach that level. The polynomial-kernel GG-KM
similarly deteriorates on the periodic surface. In contrast,
the RBF GG-KM maintains comparatively stable recovery across
the three generating mechanisms relative to the linear and
polynomial kernels, with its estimated $p_c=0.50$ contour
closely following the true level wherever it is present.
Figure~\ref{fig:cure-surface-recovery} therefore complements
the regret summary in Table~\ref{tab:simulation-regret} by
showing where prediction errors occur spatially: no single
predictor dominates across all three surfaces, with the
linear kernel performing best on the linear surface, the RBF
kernel on the quadratic surface, and the additive spline on
the periodic surface.

\GGKMWideFigure{simulation_cure_surface_recovery.pdf}{Pointwise accuracy of the three primary cure surfaces at $n=600$ under Poisson generation. Rows show the linear, quadratic and periodic generating surfaces; columns show the linear-kernel GG-KM, additive-spline, polynomial-kernel GG-KM and radial-basis-function (RBF) GG-KM predictors. The filled field is the root mean squared error of the predicted cure probability at each of the $41\times41$ grid points, taken across all valid replications, on one scale shared by every panel and truncated above 0.30. Solid labelled lines are the true cure levels $0.25$, $0.50$ and $0.60$; the dashed line is the mean estimated $0.50$ level and is annotated as absent when the fitted range does not contain it. The value printed in each panel is the mean of the displayed field. Failed fits remain in $n_{\mathrm{total}}$ and reduce the pointwise $n_{\mathrm{valid}}$ rather than disappearing.}{fig:cure-surface-recovery}{A three-row by four-column grid of shaded panels shows the pointwise root mean squared error of the predicted cure probability for the linear, quadratic and periodic generating surfaces. Columns correspond to the linear-kernel GG-KM, additive-spline, polynomial-kernel GG-KM and RBF GG-KM predictors. Darker shading marks larger error, and the printed value in each panel is the mean of the displayed field. Solid labelled lines mark the true cure levels, while a dashed line marks the mean estimated 0.50 level and is explicitly reported as absent wherever that level is outside the fitted range.}

The complete mean IPCW integrated Brier scores for all eight methods,
with their Monte Carlo standard errors, appear in the Supporting
Information. They confirm the geometry-specific ordering in
Table~\ref{tab:simulation-best} and retain both fitted count laws for
every predictor.

To complement the scenario-specific comparisons, we
summarized performance using regret relative to the
best-performing method within the retained comparison set.
Table~\ref{tab:simulation-regret} restricts this analysis to
the Poisson and negative-binomial spline models and RBF GG-KM
models across the 18 scenarios. For IPCW integrated Brier score,
the Poisson spline had the smallest maximum paired regret,
0.00065 with a 95\% Monte Carlo half-width of 0.00056, whereas
the negative-binomial spline had the smallest equally weighted
mean regret, 0.00012 (half-width 0.00010). For cure root mean
squared error, the Poisson spline had the smallest maximum regret,
0.00464 (half-width 0.00210), whereas the negative-binomial spline
had the smallest mean regret, 0.00073 (half-width 0.00025). In
contrast, the negative-binomial RBF GG-KM was most competitive
for cure AUC, attaining the smallest maximum regret, 0.00687
(half-width 0.00464), the smallest mean regret, 0.00144
(half-width 0.00060), and the best descriptive mean rank of
2.06. Thus, the regret summaries favor the spline models for
prediction error but show greater competitiveness of the RBF
GG-KM for cure discrimination.

\GGKMTable{sections/generated/simulation-regret-table}

\subsection{Higher-dimensional interaction surfaces}\label{sec:highdim-simulation}

The four-covariate experiment is a secondary mechanism study developed
after the primary analysis. It probes interaction structures that
an additive spline cannot represent, but it is neither
independent confirmation nor evidence of a general kernel
advantage. Covariates satisfy
$x_1,\dots,x_4\sim\operatorname{Uniform}(-1,1)$, and the first
target is the \emph{interaction} surface
\begin{equation}
p_c(\bm{x})=\operatorname{expit}\bigl(C_{\mathrm{int}}+1.60\,x_1x_2-1.30\,x_3x_4+0.70\,x_1-0.50\,x_4\bigr),
\end{equation}
and the \emph{radial} surface is
\begin{equation}
p_c(\bm{x})=\operatorname{expit}\bigl(C_{\mathrm{rad}}+3.00\exp\{-0.80(x_1^2+x_2^2+x_3^2+x_4^2)\}\bigr).
\end{equation}
Intercepts $C_{\mathrm{int}}$ and $C_{\mathrm{rad}}$ were
calibrated by bisection on a fixed sample of $200\,000$
observations (seed 20260901) so that each surface yields a mean
cure fraction of $0.45$. The latency remained the common
generalized gamma with $(a,d,p)=(1.6,1.3,1.1)$ and no
covariate-dependent scale; censoring was exponential with mean
3.5. The same fitted methods as in the primary experiment were
compared under Poisson, negative-binomial ($\phi=0.25$), and
negative-binomial ($\phi=1.00$) generating laws at $n=300$ and
$n=600$, giving 12 scenarios. The final experiment used 250
replicates per scenario, and every declared cell remains in the
analysis regardless of its ordering.

Table~\ref{tab:highdim-best} and
Figure~\ref{fig:highdim-paired-differences} report the directed
comparisons with paired Monte Carlo uncertainty. Their
interpretation is limited to the evaluated interaction and
radial mechanisms, and it retains every declared cell,
including comparisons that do not favor GG-KM.

    \GGKMTable{sections/generated/highdim-best-table}
\GGKMWideFigure
{highdim_rbf_vs_linear_spline_polynomial_paired_differences_with_auc.pdf}
{Paired predictive differences of the radial-basis GG-KM predictor against the count-matched GG-KM linear-kernel, spline, and GG-KM polynomial-kernel comparators on the interaction and radial cure surfaces. The top row shows paired mean IPCW integrated Brier score differences and the bottom row paired mean cure root mean squared error differences. Each plotted contrast is defined as the RBF GG-KM value minus the corresponding comparator value, so negative differences favor the RBF GG-KM and positive differences favor the comparator. Solid lines with filled markers denote Poisson fits, whereas dashed lines with open markers denote negative-binomial fits. Each panel displays the six generating-law and sample-size cells, with error bars spanning $\pm1.96$ paired Monte Carlo standard errors.}
{fig:highdim-paired-differences}
{Two-row by two-column grid of paired-difference plots for the interaction and radial cure surfaces, comparing RBF GG-KM with GG-KM linear-kernel, spline, and GG-KM polynomial-kernel models under Poisson and negative-binomial latent-count specifications.}

The complete per-scenario means with Monte Carlo standard errors, together with the numerical diagnostics, are
reported in the Supporting Information.
